\documentclass{article}

\usepackage[utf8]{inputenc}
\usepackage[T1]{fontenc}
\usepackage{helvet}
\usepackage{geometry}
\usepackage{xcolor}
\usepackage{eso-pic}
\usepackage{fancyhdr}
\usepackage{parskip}
\usepackage{microtype}
\usepackage[english, ukrainian]{babel}
\usepackage{url}
\usepackage{amsmath}
\usepackage{amssymb}
\usepackage{amsthm}
\usepackage{longtable}
\usepackage{booktabs}
\usepackage{hyperref}
\usepackage{titlesec}

\definecolor{nato}{HTML}{293757}
\definecolor{footergray}{gray}{0.65}

\geometry{a4paper,left=3cm,right=3cm,top=3cm,bottom=3cm}
\renewcommand{\familydefault}{\sfdefault}
\setcounter{secnumdepth}{3}

\titleformat{\section}
  {\color{nato}\Huge\bfseries}{\thesection.}{1em}{}
\titlespacing*{\section}{0pt}{30pt}{12pt}

\titleformat{\subsection}
  {\color{nato}\LARGE\bfseries}{\thesubsection.}{1em}{}
\titlespacing*{\subsection}{0pt}{20pt}{8pt}

\titleformat{\subsubsection}
  {\color{nato}\Large\bfseries}{\thesubsubsection.}{1em}{}
\titlespacing*{\subsubsection}{0pt}{10pt}{6pt}

\fancyhf{}
\fancyhead[R]{\small\color{footergray} 31 березня 2026}
\fancyfoot[L]{\small\color{footergray} Технічні вимоги ЄСІКС. IV. Реєстри}

\fancyfoot[R]{\small\color{footergray}\thepage}

\newcommand\HUGE[1]{\fontsize{40}{40}\selectfont #1}

\renewcommand{\headrulewidth}{0pt}
\renewcommand{\footrulewidth}{0.4pt}
\renewcommand{\footrule}{\color{footergray}\hrule width \textwidth height 0.4pt}

\begin{document}

\pagestyle{fancy}

\begin{titlepage}
  \AddToShipoutPictureBG*{\AtPageLowerLeft{\color{nato}\rule{\paperwidth}{\paperheight}}}
  \color{white}\thispagestyle{empty}\vspace*{4cm}\centering
  {\Huge  \textbf{Технічні вимоги ЄСІКС} \par} \vspace{2cm}
  {\HUGE  \textbf{IV. Реєстри} \par} \vspace{2.5cm}
  {\LARGE \textbf{IV-1. Виконавчі документи} \par}
  {\LARGE \textbf{IV-2. Судові рішення} \par}
  {\LARGE \textbf{IV-3. Правові позиції} \par}
  {\LARGE \textbf{IV-4. Дайджести} \par}
  {\LARGE \textbf{IV-5. Фізичні особи} \par}
  {\LARGE \textbf{IV-6. Адвокати} \par}
  {\LARGE \textbf{IV-7. Уповноважені} \par}
  {\LARGE \textbf{IV-8. Кадри судової влади} \par}
  {\LARGE \textbf{IV-9. Суб'єкти господарювання} \par}
  {\LARGE \textbf{IV-10. Зовнішні реєстри (ШБО)} \par}
  \vfill
  {\large 2026 \copyright\ Інформаційні Судові Системи \par}
\end{titlepage}

\newpage
\begin{center}
  {\Huge\color{nato}\bfseries Зміст\par}
\end{center}
\vspace{40pt}
\tableofcontents

\begin{abstract}
ТЕХНІЧНІ ВИМОГИ
НА СТВОРЕННЯ ЗАСОБУ ІНФОРМАТИЗАЦІЇ - КОМПОНЕНТІВ “РЕЄСТРИ”, МАЙСТЕР РЕЄСТРІВ, ЄДРСР, ЄДРВД В РАМКАХ
ЄДИНОЇ СУДОВОЇ ІНФОРМАЦІЙНО-КОМУНІКАЦІЙНОЇ СИСТЕМИ.
\end{abstract}

\newpage
\section*{Вступна частина}

\subsection*{Умовні скорочення та визначення}

Терміни та скорочення, що використовуються в документі:

\begin{longtable}{p{4cm}p{10cm}}
\toprule
\textbf{Термін} & \textbf{Значення} \\
\midrule
API & Автоматизовані програмні інтеграції, відомі також як “прикладний інтерфейс програмування – application programming interface” \\
БД & База Даних — Організована колекція даних, яка зберігається та управляється з використанням комп'ютерної системи; забезпечує зберігання, організацію та доступ до даних з метою ефективного використання та обробки \\
БП & Бізнес Процес — Сукупність дій необхідних для здійснення людиною або системою, в результаті яких відбувається досягнення бажаного результату \\
КЕП & Кваліфікований електронний підпис (X.509) \\
ЄСІКС & Єдина судова інформаційно-комунікаційна система \\
\bottomrule
\end{longtable}

\subsection*{Загальні положення}

Технічні вимоги розроблені на основі «Концепції ЄСІКС» від 2024 року \cite{esiks}, API судової системи ЄДРСР \cite{eCourt} та відповідно до Постанови Кабінету Міністрів України № 205 від 21.02.2025 \cite{law:kmu205}. Формалізація та структура технічних вимог до компонентів реєстрів виконана відповідно до міжнародного стандарту ISO/IEC/IEEE 29148:2018 \cite{iso:29148}.

\newpage
\section{Реєстр виконавчих документів}

\subsection{Вимоги до сутностей}

\subsubsection{Виконавчі листи}
\subsubsection{Судові накази}
\subsubsection{Виконавчі написи нотаріусів}
\subsubsection{Ухвали}
\subsubsection{Посвідчення комісій по трудових спорах}
\subsubsection{Постанови органів}
\subsubsection{Постанови державних виконавців}
\subsubsection{Постанови приватних виконавців}
\subsubsection{Рішення і акти виконавчих органів}
\subsubsection{Рішення Європейського суду з прав людини}
\subsubsection{Рішення (постанови) суб’єктів державного фінансового моніторингу}
\subsubsection{Рішення Органу суспільного нагляду за аудиторською діяльністю або Аудиторської палати України}
\subsubsection{Рішення Національної ради України з питань телебачення і радіомовлення}

\subsection{Вимоги до процесів}

\subsubsection{Публікація наказу}
\subsubsection{Публікація ухвали}
\subsubsection{Публікація постанови}
\subsubsection{Публікація рішення}
\subsubsection{Життєвий цикл посвічення}
\subsubsection{Життєвий цикл виконавчого напису}

\newpage
\section{Реєстр судових рішень}
Технічні вимоги до реєстру судових рішень в повній мірі реалізуються ЄДРСР (версія 1.16) та включають розробку і функціонування фронт-фасаду (інтеграційного шлюзу API), який забезпечує індексний доступ до бази даних та отримання відповідних PDF-файлів з об'єктного сховища.

\subsection{Архітектура фронт-фасаду та системна інтеграція}

Для забезпечення високої швидкодії, безпеки та відмовостійкості взаємодії компонентів реєстру встановлюються такі архітектурні вимоги:

\begin{enumerate}
    \item \textbf{Розділення рівнів (Separation of Concerns):} Веб-клієнт реєстру (Angular Frontend) не взаємодіє напряму із внутрішніми базами даних та сховищами файлів. Усі запити проходять через проміжний Фронт-Фасад (Front Facade API Gateway / Integration Service).
    \item \textbf{Індексний та повнотекстовий доступ (InterSystems Caché \& Elasticsearch):} 
    \begin{itemize}
        \item Доступ до метаданих та індексів судових рішень здійснюється через СУБД InterSystems Caché, а повнотекстовий пошук — через інтегрований Elasticsearch.
        \item Фронт-фасад транслює пошукові запити користувачів у параметризовані SQL-запити до таблиць індексів Caché через TCP-з'єднання (ODBC/JDBC або REST API).
        \item Пошук здійснюється за обов'язковими атрибутами: найменування/код суду, форма судочинства, стадія розгляду, форма рішення, дата ухвалення, 17-значний номер кримінального провадження, унікальний ідентифікатор судового рішення, ЄУН справи, склад суду (ПІБ суддів), а також ПІБ учасників справи (лише для авторизованих користувачів з повним доступом).
        \item Пошукова система підтримує предиктивний пошук (автозаповнення на основі історії та популярних запитів), збереження пошукових схем для автоматичного відстеження та алгоритми на основі теорії графів для пошуку зв'язків між пов'язаними справами та ухвалами.
    \end{itemize}
    \item \textbf{Зберігання та отримання документів (Scality Object Storage):}
    \begin{itemize}
        \item Документи (оригінали та образи) зберігаються у форматах PDF, RTF/DOCX та HTML в об'єктному сховищі Scality (з використанням S3-сумісного інтерфейсу або CDMI).
        \item Доступ до файлів виконується за унікальними ключами об'єктів (UUID/OID), які зберігаються в індексній базі даних Caché.
    \end{itemize}
    \item \textbf{Деперсоналізація та безпека (Security \& Role-Based Access):}
    \begin{itemize}
        \item Фронт-фасад валідує авторизаційні JWT-токени та ролі користувачів з інтеграцією через Електронний кабінет.
        \item Автоматичне знеособлення (деперсоналізація) текстів виконується шляхом першочергової інтеграції з діючим сервісом «Акула» та використанням ШІ (NLP) для видалення адрес, телефонів, email, РНОКПП, реквізитів документів, УНЗР та імен суддів у кримінальних провадженнях.
        \item Для авторизованих адміністраторів підтримується двовіконний режим порівняння повного та знеособленого текстів судового рішення side-by-side з можливістю ручної верифікації та коригування.
        \item Гіперпосилання на судові рішення формуються за алгоритмом криптографічного шифрування ідентифікатора документа (`CryptoService`) для запобігання перебору (prediction/enumeration) та несанкціонованого вивантаження даних.
        \item Кожне завантаження PDF-документа логується на стороні Фронт-фасаду в журналі аудиту. Максимальний час відповіді сервісів не повинен перевищувати 2 секунди.
    \end{itemize}
\end{enumerate}


\subsection{Вимоги до довідників}

\subsubsection{Категорія справи (CaseCategory)}
Довідник категорій судових справ визначає класифікацію справ за суттю спору та кодами правовідносин.
\begin{itemize}
    \item \texttt{CscId} (Ціле число) -- унікальний ідентифікатор категорії справи.
    \item \texttt{Code} (Рядок) -- цифровий/символьний код категорії.
    \item \texttt{Name} (Рядок) -- повна назва категорії судової справи.
    \item \texttt{Csc\_Csc\_Id} (Ціле число) -- посилання на батьківський елемент ієрархії категорій.
    \item \texttt{SearchCode} (Рядок) -- спеціалізований пошуковий індексний код.
\end{itemize}

\subsubsection{Роль судді у справі / провадженні (CaseJudgeRole)}
Довідник ролей суддів визначає статус та повноваження судді в складі колегії.
\begin{itemize}
    \item \texttt{RoleId} (Ціле число) -- унікальний ідентифікатор ролі.
    \item \texttt{Name} (Рядок) -- назва ролі (наприклад: Головуючий суддя, Суддя-доповідач, Член колегії).
    \item \texttt{Code} (Рядок) -- системний код ролі судді.
\end{itemize}

\subsubsection{Тип (роль) учасника справи (CaseMemberRole)}
Довідник ролей учасників справи (процесуального статусу сторін).
\begin{itemize}
    \item \texttt{Id} (Ціле число) -- унікальний ідентифікатор статусу.
    \item \texttt{Name} (Рядок) -- процесуальний статус (наприклад: Позивач, Відповідач, Третя особа, Обвинувачений, Захисник, Потерпілий).
\end{itemize}

\subsubsection{Стан розгляду справи (CaseStatus)}
Довідник станів розгляду судової справи.
\begin{itemize}
    \item \texttt{StatusId} (Ціле число) -- унікальний ідентифікатор стану.
    \item \texttt{Name} (Рядок) -- найменування стану розгляду (наприклад: Зареєстровано, Призначено до розгляду, Зупинено провадження, Завершено розгляд).
\end{itemize}

\subsubsection{Суд (Court)}
Довідник судових установ України.
\begin{itemize}
    \item \texttt{CrtId} (Ціле число) -- унікальний ідентифікатор суду.
    \item \texttt{Code} (Рядок) -- 4-значний унікальний код суду згідно з класифікатором.
    \item \texttt{Name} (Рядок) -- повне офіційне найменування суду.
    \item \texttt{RcCrtType} (Ціле число) -- посилання на тип суду.
\end{itemize}

\subsubsection{Тип суду (CourtType)}
Довідник типів судових органів.
\begin{itemize}
    \item \texttt{TypeId} (Ціле число) -- унікальний ідентифікатор типу суду.
    \item \texttt{Name} (Рядок) -- назва типу суду (наприклад: Місцевий загальний, Апеляційний, Касаційний, Верховний Суд).
\end{itemize}

\subsubsection{Категорія документу (DocCategory)}
Довідник категорій документів (судових чи супровідних).
\begin{itemize}
    \item \texttt{CategoryId} (Ціле число) -- унікальний ідентифікатор категорії.
    \item \texttt{Name} (Рядок) -- опис категорії (наприклад: Судові рішення, Супровідні документи, Технічні матеріали, Службове листування).
\end{itemize}

\subsubsection{Тип документу (DocType)}
Довідник типів судових рішень та процесуальних документів.
\begin{itemize}
    \item \texttt{Id} (Ціле число) -- унікальний ідентифікатор типу.
    \item \texttt{Name} (Рядок) -- найменування типу документа (наприклад: Вирок, Рішення, Постанова, Ухвала, Окрема думка судді).
    \item \texttt{Code} (Рядок) -- системний буквений або цифровий код типу.
\end{itemize}

\subsubsection{Юрисдикція суду (JurisdictionType)}
Довідник видів юрисдикції/спеціалізації судів.
\begin{itemize}
    \item \texttt{Id} (Ціле число) -- унікальний ідентифікатор юрисдикції.
    \item \texttt{Name} (Рядок) -- опис спеціалізації суду (наприклад: Загальна, Господарська, Адміністративна, Військова).
\end{itemize}

\subsubsection{Вид судочинства (JusticeType)}
Довідник видів судочинства.
\begin{itemize}
    \item \texttt{Id} (Ціле число) -- унікальний ідентифікатор виду судочинства.
    \item \texttt{Name} (Рядок) -- назва (наприклад: Цивільне, Кримінальне, Господарське, Адміністративне, Конституційне, Справи про адміністративні правопорушення).
\end{itemize}

\subsubsection{Результат розгляду провадження (ProcResult)}
Довідник результатів вирішення судових спорів та проваджень.
\begin{itemize}
    \item \texttt{ResultId} (Ціле число) -- унікальний ідентифікатор результату.
    \item \texttt{Name} (Рядок) -- опис результату розгляду (наприклад: Позов задоволено повністю, Позов задоволено частково, У задоволенні відмовлено, Провадження закрито).
\end{itemize}

\subsubsection{Стадія розгляду провадження (ProcStage)}
Довідник процесуальних стадій розгляду судової справи.
\begin{itemize}
    \item \texttt{StageId} (Ціле число) -- унікальний ідентифікатор стадії.
    \item \texttt{Name} (Рядок) -- опис процесуальної стадії (наприклад: Перша інстанція, Апеляційний перегляд, Касаційний перегляд, Перегляд за нововиявленими обставинами).
\end{itemize}

\subsubsection{Стан розгляду провадження (ProcStatus)}
Довідник станів розгляду судового провадження.
\begin{itemize}
    \item \texttt{StatusId} (Ціле число) -- унікальний ідентифікатор стану.
    \item \texttt{Name} (Рядок) -- опис поточного статусу розгляду провадження.
\end{itemize}

\subsubsection{Регіон (Region)}
Довідник адміністративно-територіальних регіонів України.
\begin{itemize}
    \item \texttt{RegId} (Ціле число) -- унікальний ідентифікатор регіону (області).
    \item \texttt{Name} (Рядок) -- офіційне найменування області чи АР Крим.
    \item \texttt{Code} (Рядок) -- цифровий або літерний код регіону.
\end{itemize}

\subsubsection{Тип заяви (ClaimType)}
Довідник видів процесуальних заяв та позовів.
\begin{itemize}
    \item \texttt{TypeId} (Ціле число) -- унікальний ідентифікатор типу заяви.
    \item \texttt{Name} (Рядок) -- назва заяви (наприклад: Позовна заява, Заява про забезпечення позову, Заява про перегляд рішення).
\end{itemize}

\addtocontents{toc}{\protect\newpage}
\newpage
\subsection{Вимоги до сутностей}

\subsubsection{Судові рішення (DocumentDto)}
Ключова сутність, що містить основні метадані та посилання на файли судових рішень.
\begin{itemize}
    \item \texttt{documentID} (Ціле число) -- унікальний ідентифікатор судового рішення в системі.
    \item \texttt{docRegNum} (Рядок) -- державний реєстраційний номер судового рішення.
    \item \texttt{docDateReg} (Дата) -- дата реєстрації документа в реєстрі.
    \item \texttt{docDatePublish} (Дата) -- дата офіційного оприлюднення документа.
    \item \texttt{docDateSend} (Дата) -- дата надсилання документа до ЄДРСР з суду.
    \item \texttt{proceedingNumber} (Рядок) -- унікальний номер судового провадження.
    \item \texttt{criminalNumber} (Рядок) -- 17-значний номер кримінального провадження (за наявності).
    \item \texttt{crtCode} (Рядок) -- код суду, який ухвалив судове рішення.
    \item \texttt{is\_Public} (Логічний) -- ознака дозволу на публічний доступ.
    \item \texttt{cscText} (Текст) -- повний текст судового рішення (у форматі HTML або RTF).
\end{itemize}

\subsubsection{Файли документу (DocumentFileDto)}
Сутність, що містить посилання на фізичні файли документів в об'єктному сховищі.
\begin{itemize}
    \item \texttt{FileId} (Ціле число) -- унікальний ідентифікатор запису про файл.
    \item \texttt{DocumentId} (Ціле число) -- посилання на ідентифікатор судового рішення.
    \item \texttt{ScalityOriginalKey} (Рядок) -- унікальний UUID/OID оригінального PDF-файлу з КЕП у сховищі Scality.
    \item \texttt{ScalityImpersonalKey} (Рядок) -- унікальний UUID/OID знеособленого образу PDF-файлу в Scality.
    \item \texttt{SrcTypeImage} (Ціле число) -- формат вихідного файлу (PDF, RTF, HTML, PNG).
\end{itemize}

\subsubsection{Ідентифікаційні документи (DocumentAsdsLink)}
Зв'язок між рішенням у ЄДРСР та первинною справою в автоматизованій системі документообігу суду (АСДС).
\begin{itemize}
    \item \texttt{LinkId} (Ціле число) -- унікальний ідентифікатор зв'язку.
    \item \texttt{DocumentId} (Ціле число) -- посилання на судове рішення.
    \item \texttt{AsdsSystemId} (Ціле число) -- ідентифікатор версії/екземпляра АСДС.
    \item \texttt{AsdsDocumentId} (Рядок) -- унікальний внутрішній номер документа в АСДС.
\end{itemize}

\subsubsection{Заява (ClaimPostDto)}
Сутність електронного позову чи заяви, що направляється через Електронний суд.
\begin{itemize}
    \item \texttt{ClaimId} (Ціле число) -- унікальний ідентифікатор заяви.
    \item \texttt{CaseNumber} (Рядок) -- унікальний номер справи, до якої подається заява.
    \item \texttt{ClaimType} (Ціле число) -- посилання на тип заяви.
    \item \texttt{SenderId} (Ціле число) -- ідентифікатор особи-відправника.
    \item \texttt{SubmissionDate} (Дата і час) -- дата та час відправлення.
\end{itemize}

\subsubsection{Квинтація заяви (ClaimTicketDto)}
Квитанція про реєстрацію заяви в системі документообігу.
\begin{itemize}
    \item \texttt{TicketId} (Ціле число) -- унікальний ідентифікатор квитанції.
    \item \texttt{ClaimId} (Ціле число) -- посилання на відповідну заяву.
    \item \texttt{TicketNumber} (Рядок) -- системний реєстраційний номер квитанції.
    \item \texttt{Status} (Ціле число) -- статус обробки заяви (наприклад: Зареєстровано, Відхилено).
    \item \texttt{LogDate} (Дата і час) -- час формування квитанції.
\end{itemize}

\subsubsection{Справа (Case)}
Судова справа, що містить сукупність проваджень та документів.
\begin{itemize}
    \item \texttt{CaseId} (Ціле число) -- унікальний ідентифікатор судової справи.
    \item \texttt{CaseNumber} (Рядок) -- єдиний унікальний номер судової справи (ЄУН).
    \item \texttt{RegistrationDate} (Дата) -- дата первинної реєстрації справи в суді.
    \item \texttt{CategoryCode} (Рядок) -- посилання на категорію справи.
\end{itemize}

\subsubsection{Провадження (Proceedings)}
Конкретна стадія або епізод розгляду судової справи.
\begin{itemize}
    \item \texttt{ProcId} (Ціле число) -- унікальний ідентифікатор провадження.
    \item \texttt{CaseId} (Ціле число) -- посилання на судову справу.
    \item \texttt{ProcNumber} (Рядок) -- номер провадження в суді.
    \item \texttt{Stage} (Ціле число) -- посилання на стадію розгляду.
\end{itemize}

\subsubsection{Учасник справи (CaseMemberDto)}
Фізична або юридична особа, яка є стороною процесу.
\begin{itemize}
    \item \texttt{MemberId} (Ціле число) -- унікальний ідентифікатор учасника.
    \item \texttt{CaseId} (Ціле число) -- посилання на судову справу.
    \item \texttt{FullName} (Рядок) -- ПІБ фізичної особи або назва юридичної особи (зашифровано/знеособлено для публічного доступу).
    \item \texttt{Role} (Ціле число) -- посилання на процесуальну роль учасника.
\end{itemize}

\subsubsection{Суддя по справі (CaseJudge)}
Дані про суддю або колегію суддів, які розглядають справу.
\begin{itemize}
    \item \texttt{Id} (Ціле число) -- унікальний ідентифікатор призначення судді.
    \item \texttt{CaseId} (Ціле число) -- посилання на судову справу.
    \item \texttt{JudgeFio} (Рядок) -- повне ім'я (ПІБ) судді.
    \item \texttt{Role} (Ціле число) -- посилання на роль судді в провадженні.
\end{itemize}

\subsubsection{Календар судових засідань (Trial)}
Дані про призначені судові засідання за справою.
\begin{itemize}
    \item \texttt{TrialId} (Ціле число) -- унікальний ідентифікатор засідання.
    \item \texttt{CaseId} (Ціле число) -- посилання на судову справу.
    \item \texttt{TrialDate} (Дата і час) -- дата та час проведення судового засідання.
    \item \texttt{CourtRoom} (Рядок) -- номер або назва залу судового засідання.
\end{itemize}

\subsubsection{Виконавче провадження (ExecutiveDecision)}
Реквізити виконавчого документа, виданого за результатом розгляду справи.
\begin{itemize}
    \item \texttt{DecisionId} (Ціле число) -- унікальний ідентифікатор рішення.
    \item \texttt{DocumentId} (Ціле число) -- посилання на судове рішення, на підставі якого видано документ.
    \item \texttt{ExecutiveNumber} (Рядок) -- реєстраційний номер виконавчого провадження.
    \item \texttt{State} (Ціле число) -- статус виконання.
\end{itemize}

\subsubsection{Довіреність (PowerOfAttorney)}
Дані про електронні довіреності представників сторін у справі.
\begin{itemize}
    \item \texttt{Id} (Ціле число) -- унікальний ідентифікатор довіреності.
    \item \texttt{PrincipalId} (Ціле число) -- ідентифікатор довірителя.
    \item \texttt{RepresentativeId} (Ціле число) -- ідентифікатор представника (адвоката).
    \item \texttt{ExpiryDate} (Дата) -- термін дії довіреності.
\end{itemize}

\subsubsection{Документи сторін (PartyDocPostDto)}
Електронні документи та докази, що подаються сторонами процесу.
\begin{itemize}
    \item \texttt{Id} (Ціле число) -- унікальний ідентифікатор документа.
    \item \texttt{ClaimId} (Ціле число) -- посилання на заяву, до якої додається документ.
    \item \texttt{FileName} (Рядок) -- назва файлу.
    \item \texttt{PayloadKey} (Рядок) -- ключ файлу в об'єктному сховищі Scality.
\end{itemize}

\subsubsection{Постанова виконавчого провадження (ExecutiveDecisionDocumentDto)}
Постанови та документи державних або приватних виконавців.
\begin{itemize}
    \item \texttt{Id} (Ціле число) -- унікальний ідентифікатор постанови.
    \item \texttt{DecisionId} (Ціле число) -- посилання на виконавче провадження.
    \item \texttt{DocNumber} (Рядок) -- реєстраційний номер постанови.
    \item \texttt{DocType} (Ціле число) -- тип постанови виконавця.
\end{itemize}

\subsubsection{Результат розсилки документів сторонам (PartyDocPostResultDto)}
Статус розсилки копій судових документів учасникам справи.
\begin{itemize}
    \item \texttt{Id} (Ціле число) -- унікальний ідентифікатор результату розсилки.
    \item \texttt{DocumentId} (Ціле число) -- посилання на розіслане судове рішення.
    \item \texttt{RecipientId} (Ціле число) -- ідентифікатор одержувача.
    \item \texttt{IsDelivered} (Логічний) -- ознака успішної доставки.
\end{itemize}

\subsubsection{Квитанція про доставку cтороні справи (CaseMemberTicketDto)}
Електронна квитанція про отримання судового рішення в кабінеті користувача.
\begin{itemize}
    \item \texttt{Id} (Ціле число) -- унікальний ідентифікатор квитанції.
    \item \texttt{DocumentId} (Ціле число) -- посилання на доставлене судове рішення.
    \item \texttt{MemberId} (Ціле число) -- посилання на учасника справи.
    \item \texttt{DeliveryDate} (Дата і час) -- точний час доставки в електронний кабінет.
\end{itemize}

\newpage
\subsection{Вимоги до процесів}

\subsubsection{Завантаження файлу}
Процес завантаження PDF-файлу судового рішення:
\begin{enumerate}
    \item Запит від клієнта на завантаження документа за його ID.
    \item Авторизація запиту та визначення ролі користувача на Фронт-фасаді.
    \item Отримання відповідного ключа файлу (оригінал/деперсоналізований) з бази даних InterSystems Caché.
    \item Запит файлу з об'єктного сховища Scality через S3/CDMI API.
    \item Передача файлу (streaming) клієнту та запис події в журнал аудиту.
\end{enumerate}

\subsubsection{Видалення файлу}
Процес обмеження доступу або логічного видалення файлу рішення:
\begin{enumerate}
    \item Запит від уповноваженого адміністратора на видалення або обмеження доступу до рішення.
    \item Перевірка КЕП та прав доступу адміністратора на Фронт-фасаді.
    \item Встановлення ознаки обмеження доступу у записі індексу СУБД InterSystems Caché.
    \item Надсилання запиту до Scality на видалення або архівацію об'єкта PDF.
    \item Оновлення пошукових індексів для виключення документа з публічного доступу.
    \item Фіксація події видалення в журналі аудиту.
\end{enumerate}

\subsubsection{Життєвий цикл справ}
Процес проходження судової справи в системі:
\begin{enumerate}
    \item Реєстрація справи в АСДС суду з присвоєнням єдиного унікального номера справи (ЄУН).
    \item Автоматична синхронізація метаданих справи з центральним реєстром.
    \item Формування проваджень на різних стадіях розгляду (перша, апеляційна, касаційна інстанції).
    \item Прив'язка ухвалених судових рішень та призначених засідань до провадження справи.
    \item Закриття справи після набрання остаточним судовим рішенням законної сили та перенесення метаданих в архівний стан.
\end{enumerate}

\subsubsection{Життєвий цикл заяви}
Процес подачі та реєстрації електронної заяви:
\begin{enumerate}
    \item Створення заяви та додатків користувачем в Електронному кабінеті.
    \item Підписання заяви КЕП заявника та відправка до черги.
    \item Реєстрація заяви в АСДС суду та призначення судді-доповідача.
    \item Формування та відправка квитанції про реєстрацію із зазначенням реквізитів справи та судді.
\end{enumerate}

\subsubsection{Життєвий цикл довіреності}
Процес створення та валідації електронних довіреностей:
\begin{enumerate}
    \item Створення довіреності довірителем в Електронному кабінеті із зазначенням представника та обсягу прав.
    \item Підписання довіреності КЕП довірителя та її реєстрація в реєстрі довіреностей.
    \item Перевірка довіреності Фронт-фасадом під час спроби доступу представника до документів довірителя.
    \item Анулювання довіреності після закінчення терміну дії або її дострокового відкликання.
\end{enumerate}

\subsubsection{Життєвий цикл виконавчого провадження}
Процес виконання судового рішення:
\begin{enumerate}
    \item Формування виконавчого документа в суді та його підписання КЕП судді.
    \item Автоматичне надсилання виконавчого документа до реєстру виконавчих документів (ЄДРВД).
    \item Відкриття виконавчого провадження державним або приватним виконавцем.
    \item Внесення постанов виконавця та оновлення статусу виконання в реєстрі.
    \item Закриття провадження за результатом виконання.
\end{enumerate}

\subsubsection{Розсилка документів сторонам}
Процес автоматичного надсилання документів учасникам справи:
\begin{enumerate}
    \item Реєстрація та підписання судового рішення КЕП в судовому реєстрі.
    \item Визначення переліку учасників судового процесу.
    \item Перевірка наявності Електронних кабінетів у учасників справи.
    \item Автоматична розсилка копій рішень в Електронні кабінети сторін.
    \item Формування завдання на паперову розсилку для учасників без Електронного кабінету.
\end{enumerate}

\subsubsection{Отримання квитанції про розсилку сторонам}
Процес фіксації отримання документів сторонами:
\begin{enumerate}
    \item Доставка документа в Електронний кабінет сторони справи.
    \item Генерація квитанції про доставку сервером кабінету з фіксацією дати та часу (time-stamp).
    \item Підписання квитанції КЕП системи та її збереження в реєстрі індексів.
    \item Відображення квитанції у справі як підтвердження належного сповіщення сторони.
\end{enumerate}

\subsubsection{Подача заяви в чергу на обробку}
Процес асинхронного надсилання заяви до суду:
\begin{enumerate}
    \item Надсилання підписаної заяви з додатками через API Фронт-фасаду.
    \item Валідація КЕП та поміщення пакету документів у чергу повідомлень (RabbitMQ/Kafka).
    \item Повернення клієнту статусу прийняття в чергу та унікального ID транзакції.
    \item Зчитування пакета воркером та передача його в АСДС відповідного суду.
\end{enumerate}

\subsubsection{Отримання квитанції про статус заяви}
Процес контролю реєстрації заяви:
\begin{enumerate}
    \item Запит від веб-клієнта щодо статусу обробки заяви за ID транзакції.
    \item Отримання від АСДС суду квитанції про реєстрацію заяви з присвоєним номером.
    \item Збереження квитанції в базі даних та оновлення статусу заяви.
    \item Повернення підписаної квитанції про успішну реєстрацію клієнту.
\end{enumerate}
\subsubsection{Публікація судового рішення}
Процес публікації нового судового рішення:
\begin{enumerate}
    \item Передача оригінального електронного примірника рішення з КЕП суддів до ЄДРСР.
    \item Автоматичне знеособлення документа за допомогою інтеграції з сервісом «Акула» (або лінгвістичним ШІ-сервісом) для вилучення відомостей обмеженого доступу та генерація деперсоналізованого образу.
    \item Завантаження оригінального та знеособленого файлів у форматах PDF, RTF/DOCX або HTML в об'єктне сховище Scality та отримання унікальних UUID ключів.
    \item Створення нового запису в реєстрі індексів СУБД InterSystems Caché з пошуковими метаданими та посиланнями на UUID ключів у Scality.
    \item Генерація криптографічно захищеного непередбачуваного гіперпосилання на публікацію для веб-клієнта.
    \item Оновлення індексів пошукової бази для забезпечення миттєвого доступу до опублікованого рішення.
\end{enumerate}

\subsubsection{Перевірка наявності кабінету особи}
Процес перевірки реєстрації особи в системі Електронного суду:
\begin{enumerate}
    \item Запит від АСДС або Фронт-фасаду щодо наявності кабінету особи за її РНОКПП або кодом ЄДРПОУ.
    \item Пошук відповідного профілю користувача в реєстрі кабінетів.
    \item Повернення результату перевірки (наявність кабінету, ID профілю).
    \item Вибір відповідного каналу доставки судового рішення стороні процесу.
\end{enumerate}

\newpage
\section{Реєстр правових позицій}

\subsection{Вимоги до сутностей}

\subsubsection{Правова позиція}

\subsection{Вимоги до процесів}

\subsubsection{Життєвий цикл правової позиції}

\newpage
\section{Реєстр дайджестів}

\subsection{Вимоги до сутностей}

\subsubsection{Дайджест}

\subsection{Вимоги до процесів}

\subsubsection{Життєвий цикл дайджеста}

\addtocontents{toc}{\protect\newpage}
\newpage
\section{Реєстр фізичних осіб}

\subsection{Вимоги до сутностей}

\subsubsection{Фізична особа}

\subsection{Вимоги до процесів}

\subsubsection{Життєвий цикл фізичної особи}

\newpage
\section{Реєстр адвокатів}

\subsection{Вимоги до сутностей}

\subsubsection{Адвокат}

\subsection{Вимоги до процесів}

\subsubsection{Життєвий цикл адвоката}

\newpage
\section{Реєстр суддів}

\subsection{Вимоги до сутностей}

\subsubsection{Суддя}

\subsection{Вимоги до процесів}

\subsubsection{Життєвий цикл судді}

\newpage
\section{Реєстр уповноважених}

\subsection{Вимоги до сутностей}

\subsubsection{Уповноважений}

\subsection{Вимоги до процесів}

\subsubsection{Життєвий цикл уповноваженого}

\newpage
\section{Реєстр суб'єктів Мінекономіки}

\subsection{Вимоги до сутностей}

\subsubsection{Суб'єкт господарської діяльності}

\subsection{Вимоги до процесів}

\subsubsection{Життєвий цикл суб'єкта господарювання}

\newpage
\section*{Висновок}

\begin{thebibliography}{9}

\bibitem{ecoz:rehab} Міністерство охорони здоров’я України. \newblock Технічні вимоги Реабілітація 2.0. \newblock Версія 08.08.2024.
\bibitem{dstu:3973} ДСТУ 3973-2000. \newblock Система розроблення та постачання продукції на виробництво. Правила виконання науково-дослідних робіт. Загальні положення.
\bibitem{dstu:3974} ДСТУ 3974-2000. \newblock Система розроблення та постачання продукції на виробництво. Правила виконання науково-конструкторських робіт. Загальні положення.
\bibitem{dstu:42010} ДСТУ 42010:2018. \newblock Інженерія систем і програмних засобів. Опис архітектури.
\bibitem{dstu:3008} ДСТУ 3008:2015. \newblock Інформація та документація. Звіти у сфері науки і техніки. Структура та правила оформлювання.
\bibitem{dstu:62264} ДСТУ IEC 62264-1:2019. \newblock Інтеграція систем керування підприємством та виробництвом. Частина 1. Моделі та термінологія.
\bibitem{dstu:en301549} ДСТУ EN 301 549:2022. \newblock Інформаційні технології. Вимоги щодо доступності продуктів та послуг ІКТ для людей з інвалідністю.
\bibitem{dstu:4541} ДСТУ 4541:2006. \newblock Інформаційні технології. Засоби криптографічного захисту інформації. Алгоритм симетричного шифрування.
\bibitem{iso:19510} ISO/IEC 19510:2013. \newblock Information technology -- Object Management Group Business Process Model and Notation (BPMN).
\bibitem{law:info-protection} Закон України «Про захист інформації в інформаційно-телекомунікаційних системах».
\bibitem{law:personal-data} Закон України «Про захист персональних даних».
\bibitem{law:cybersecurity} Закон України «Про основні засади забезпечення кібербезпеки України».
\bibitem{owasp:top10} OWASP Foundation. \newblock OWASP Top 10: The Ten Most Critical Web Application Security Risks. \newblock \url{https://owasp.org/www-project-top-ten/}.
\bibitem{nist:sp800-57} NIST Special Publication 800-57 Part 1 Revision 5. \newblock Recommendation for Key Management.
\bibitem{iso:42010} ISO/IEC/IEEE 42010:2022. \newblock Systems and software engineering — Architecture description.
\bibitem{zachman} Zachman, J.A. \newblock A Framework for Information Systems Architecture. \newblock IBM Systems Journal, 1987.
\bibitem{esiks} Концепція ЄСІКС. \url{https://court.gov.ua/storage/portal/dsa/news/Kontseptsia%20ESIKS.pdf}. 2024
\bibitem{cbor} CBOR/COSE encoding/protocol. \url{https://tonpa.guru/stream/2024/2024-11-20%20CBOR%20COSE.htm}. 2025
\bibitem{law:kmu205} Постанова Кабінету Міністрів України № 205 від 21.02.2025 «Про затвердження Порядку взаємодії та ведення державних інформаційних ресурсів в межах Єдиної судової інформаційно-комунікаційної системи».
\bibitem{iso:29148} ISO/IEC/IEEE 29148:2018. \newblock Systems and software engineering — Life cycle processes — Requirements engineering.

\end{thebibliography}

\end{document}
